%===============================================================================
% LaTeX-sjabloon voor het rapport van het graduaatsproject
% Graduaat in het Programmeren - HOGENT
%
% Hoofddocument van het graduaatsproject.
% De persoonlijke gegevens en projectgegevens worden beheerd in
% ../projectgegevens.tex. De inhoud van de verschillende hoofdstukken
% staat in afzonderlijke .tex-bestanden.
%===============================================================================
TODO! LEES ALLES NA 
\documentclass[dutch,dit,report]{hogentreport}

%===============================================================================
% Algemene instellingen
%===============================================================================
%% Afbeeldingen komen in de map graphics/
\graphicspath{{../graphics/}}
%===============================================================================
% Codefragmenten
%===============================================================================
% Codefragmenten worden weergegeven met minted.
% Hiervoor is compilatie met -shell-escape vereist.
%
% Bij gebruik van een oudere TeX Live-installatie (TeX Live 2022 of ouder,
% met minted versie 2) kan de configuratie moeten worden aangepast volgens
% de instructies van het HOGENT-sjabloon.
\usepackage[chapter]{minted}
\usemintedstyle{solarized-light}
\setminted{%
    autogobble,
    frame=lines,
    breaklines,
    linenos,
    tabsize=4
}
% Nederlandse benamingen voor codefragmenten.
\renewcommand\listoflistingscaption{%
    \IfLanguageName{dutch}{Lijst van codefragmenten}{List of listings}
}
\renewcommand\listingscaption{%
    \IfLanguageName{dutch}{Codefragment}{Listing}
}
\renewcommand*\listoflistings{%
    \cleardoublepage\phantomsection\addcontentsline{toc}{chapter}{\listoflistingscaption}%
    \listof{listing}{\listoflistingscaption}%
}
%===============================================================================
% URL's
%===============================================================================
% Zorgt ervoor dat lange URL's correct kunnen worden afgebroken.
\usepackage{xurl}


%===============================================================================
% Projectgegevens
%===============================================================================
% Alle persoonlijke en projectgebonden gegevens worden centraal beheerd in
% projectgegevens.tex.
\input{../projectgegevens}
\author{Sarah Turner Burkitt}
\supervisor{Isabelle Demeulenaere}
\cosupervisor{Bart D'Hoore}
\title[Proof of Concept voor aanpasbare e-mailtemplates binnen Vesalius.ai]{Ontwikkeling van een visuele e-maileditor}
\academicyear{2026-2027}
\examperiod{1}          
\degreesought{\IfLanguageName{dutch}{Graduaat Programmeren}{Associate Degree in Programming}}
\institution{Vesalius.ai (dochterbedrijf binnen Endare)}


%===============================================================================
% LaTeX-afbrekingen
%===============================================================================
% Woorden waarvoor LaTeX standaard een minder goede afbreking kan kiezen.
\hyphenation{back-slash}

%===============================================================================
% Bibliografie
%===============================================================================
% Bronnen die specifiek voor het rapport werden gebruikt.
\addbibresource{bronnen.bib}
% Bronnen die reeds in het projectvoorstel werden gebruikt.
\addbibresource{../voorstel/voorstel.bib}
% Geen afzonderlijke titel voor de bibliografieheading.
\defbibheading{bibempty}{}


%===============================================================================
% Algemene opmaak
%===============================================================================
% Voorkomt lege pagina's tussen hoofdstukken.
\renewcommand{\cleardoublepage}{\clearpage}
% Iets meer verticale ruimte in tabellen voor een betere leesbaarheid.
\renewcommand{\arraystretch}{1.2}

%===============================================================================
% Document
%===============================================================================
\begin{document}


%===============================================================================
% Voorwerk
%===============================================================================
\frontmatter
% Verberg paginaverwijzingen tijdens het titelblad.
\hypersetup{pageanchor=false}
% Wanneer het rapport in het Engels wordt geschreven, wordt eerst een
% Nederlandstalig titelblad toegevoegd. Voor dit rapport wordt Nederlands
% gebruikt, waardoor dit gedeelte niet wordt uitgevoerd.
\IfLanguageName{english}{%
    \degreesought{Associate Degree In Programming}%
    \begin{otherlanguage}{dutch}%
        \maketitle%
    \end{otherlanguage}%
}{}
% Titelblad
\maketitle
% Vanaf hier worden paginaverwijzingen opnieuw geactiveerd.
\hypersetup{pageanchor=true}

%===============================================================================
% Voorwoord en samenvattingen
%===============================================================================
\input{voorwoord}
\input{samenvatting}

%===============================================================================
% Inhoudsopgave en lijsten
%===============================================================================
\tableofcontents

% Lijst van afbeeldingen
\listoffigures

% Lijst van tabellen
% Deze lijst blijft behouden omdat het rapport onder andere
% vergelijkingstabellen en testtabellen kan bevatten.
\listoftables

% Lijst van codefragmenten
\listoflistings

%===============================================================================
% Kern van het rapport
%===============================================================================
\mainmatter{}
% Richtwaarde voor de kerntekst: ongeveer 10 tot 15 pagina's.
%
% Het rapport is een technisch verslag van het graduaatsproject en geen
% uitgebreide scriptie. De nadruk ligt daarom op de gemaakte analyse,
% ontwerpkeuzes, technische realisatie, resultaten en reflectie.
\input{inleiding}
\input{technologie}
\input{ontwerp}
\input{realisatie}
\input{resultaat}
\input{reflectie}

%===============================================================================
% Bijlagen
%===============================================================================
\appendix

%-------------------------------------------------------------------------------
% Projectvoorstel
%-------------------------------------------------------------------------------
\chapter{Projectvoorstel}
Het onderwerp van dit graduaatsproject is een zelfstandige proof of concept
(PoC) van een visuele e-maileditor voor Vesalius.ai. Het projectvoorstel
beschrijft de probleemstelling, doelstellingen, functionele afbakening,
technologische keuzes en geplande validatie van de oplossing. Het voorstel
wordt afgestemd met de bedrijfsmentor en de begeleider van HOGENT.

Het oorspronkelijke projectvoorstel wordt als bijlage bij dit rapport
opgenomen. 

\input{../voorstel/voorstel-inhoud}

%-------------------------------------------------------------------------------
% Repository
%-------------------------------------------------------------------------------
\chapter{Repository}
De broncode van het graduaatsproject wordt beheerd in de interne
Bitbucket-omgeving van Vesalius.ai. Deze repository bevat de projectstructuur,
broncode en relevante configuratiebestanden die nodig zijn om de implementatie
te ontwikkelen, te onderhouden en binnen de bedrijfsomgeving te testen.

Omdat de repository deel uitmaakt van de interne ontwikkelomgeving van
Vesalius.ai, is deze niet publiek toegankelijk. De toegang wordt beheerd
volgens de interne afspraken van het bedrijf. Daarom wordt geen publieke
repositorylink in dit rapport opgenomen. Om de vertrouwelijkheid van
bedrijfsinformatie te beschermen, worden productiecredentials, API-sleutels,
vertrouwelijke configuratiegegevens en echte patiëntgegevens niet opgenomen
in het rapport. Waar relevant worden de architectuur, technische keuzes en
testresultaten toegelicht aan de hand van informatie die binnen de gemaakte
afspraken gedeeld mag worden.

\begin{center}
    \url{\repolink}
\end{center}

De repository en de daarin opgenomen broncode blijven eigendom van Vesalius.ai.
Indien bepaalde onderdelen omwille van vertrouwelijkheid, bedrijfsgevoelige
informatie of toegangsbeperkingen niet publiek beschikbaar zijn, worden deze
niet openbaar gemaakt. In het rapport worden in dat geval enkel relevante
architectuur-, test- en codefragmenten opgenomen die binnen de afspraken met
het bedrijf kunnen worden gedeeld.

%-------------------------------------------------------------------------------
% Verantwoording van AI-gebruik
%-------------------------------------------------------------------------------
\input{ai-verantwoording}

%-------------------------------------------------------------------------------
% Overige bijlagen
%-------------------------------------------------------------------------------
% Oorspronkelijk projectvoorstel
\input{bijlagen/projectvoorstel}

% Testresultaten en testmatrix
\input{bijlagen/testresultaten}

% Architectuurdiagram
\input{bijlagen/architectuur}

%===============================================================================
% Nawerk
%===============================================================================
\backmatter{}

% Compactere ruimte tussen bibliografie-items.
\setlength\bibitemsep{2pt}

% Bibliografie opnemen in de inhoudsopgave.
\printbibliography[heading=bibintoc]

\end{document}